% notebooks/notes/07_bundle_adjustment.tex
\documentclass[11pt]{article}

\input{preamble.tex}

\title{Bundle Adjustment}
\date{May 2026}

\begin{document}
\maketitle

\tableofcontents

% ====
\section{Bundle Adjustment as Nonlinear Least Squares}
% ====

Bundle adjustment jointly refines camera poses and landmark positions by minimising image-plane reprojection error.

\begin{definition}[Observation graph]
Let \( C, P \) index the camera poses and landmarks, respectively. 
\[
\mathcal C=\{1,\dots,M\},\qquad \mathcal P=\{1,\dots,N\}
\]
The observation graph \( G \) is the bipartite graph
\[
\mathcal G=(\mathcal C,\mathcal P,\mathcal O),
\qquad
\mathcal O\subseteq \mathcal C\times \mathcal P
\]
Where \((i,j)\in\mathcal O\) means that the camera in pose \(i\) observes landmark \(j\).

For each \((i,j)\in\mathcal O\), let \(z_{ij}\in\R^2\) denote the measured pixel coordinate.
\end{definition}

\begin{definition}[Reprojection residual]
Let \(T_i=(R_i,t_i)\in SE(3)\) be the world-to-camera pose of camera \(i\), and let \(X_j\in\R^3\) be landmark \(j\) in world coordinates. Define
\[
Y_{ij}\coloneqq R_iX_j+t_i\in\R^3
\]
For the pinhole projection map \(\pi:\R^3\supset\{Z>0\}\to\R^2\), the predicted pixel coordinate is
\[
\hat z_{ij}=\pi(Y_{ij})
\]
The reprojection residual is therefore
\[
r_{ij}(T_i,X_j)
\coloneqq
z_{ij}-\pi(R_iX_j+t_i)
\in\R^2
\]
\end{definition}

\begin{definition}[Bundle adjustment objective]
The bundle adjustment problem is
\[
\min_{\{T_i\},\{X_j\}}
\frac12
\sum_{(i,j)\in\mathcal O}
\|r_{ij}(T_i,X_j)\|_2^2
\]
Equivalently,
\[
\min_{\{T_i\},\{X_j\}}
\frac12
\sum_{(i,j)\in\mathcal O}
\left\|
z_{ij}-\pi(R_iX_j+t_i)
\right\|_2^2
\]
\end{definition}

\begin{proposition}[Gaussian pixel noise gives weighted reprojection least squares]
Assume independently over \((i,j)\in\mathcal O\), with \(\Sigma_{ij}\succ0\)
\[
z_{ij}=\pi(R_iX_j+t_i)+\varepsilon_{ij},
\qquad
\varepsilon_{ij}\sim\mathcal N(0,\Sigma_{ij})
\]
Then, the maximum likelihood estimate of the bundle adjustment problem minimises
\[
\frac12
\sum_{(i,j)\in\mathcal O}
r_{ij}^{\top}\Sigma_{ij}^{-1}r_{ij}
\]
In particular, if \(\Sigma_{ij}=\sigma^2I_2\), this is equivalent to ordinary reprojection least squares.
\end{proposition}

\begin{proof}
For one observation, we have
\[
p(z_{ij}\mid T_i,X_j)
=
\frac{1}{2\pi|\Sigma_{ij}|^{1/2}}
\exp\!\left(
-\frac12 r_{ij}^{\top}\Sigma_{ij}^{-1}r_{ij}
\right)
\]
Independence gives a product likelihood over \(\mathcal O\). Taking the negative logarithm and discarding constants independent of the unknowns gives
\[
\frac12
\sum_{(i,j)\in\mathcal O}
r_{ij}^{\top}\Sigma_{ij}^{-1}r_{ij}
\]
If \(\Sigma_{ij}=\sigma^2I_2\), then
\[
r_{ij}^{\top}\Sigma_{ij}^{-1}r_{ij}
=
\sigma^{-2}\|r_{ij}\|_2^2
\]
The the scalar factor \(\sigma^{-2}\) does not affect the minimiser.
\end{proof}

\begin{remark}[Information form]
Writing
\[
\Omega_{ij}\coloneqq \Sigma_{ij}^{-1}
\]
The weighted objective becomes
\[
\frac12
\sum_{(i,j)\in\mathcal O}
r_{ij}^{\top}\Omega_{ij}r_{ij}
\]
Thus, higher information corresponds to a more trusted image measurement.
\end{remark}


% ====
\section{Residual Jacobians}
% ====

We now compute the local sensitivities of a reprojection residual with respect to a landmark and a camera pose.

\begin{lemma}[Projection Jacobian]
Let
\[
Y=
\begin{bmatrix}
X_c\\Y_c\\Z_c
\end{bmatrix},
\qquad Z_c>0
\]
and
\[
\pi(Y)=
\begin{bmatrix}
f_xX_c/Z_c+c_x\\
f_yY_c/Z_c+c_y
\end{bmatrix}
\]
Then
\[
D\pi(Y)
=
\frac{\partial \pi}{\partial Y}
=
\begin{bmatrix}
f_x/Z_c & 0 & -f_xX_c/Z_c^2\\
0 & f_y/Z_c & -f_yY_c/Z_c^2
\end{bmatrix}
\]
\end{lemma}

\begin{proof}
Differentiate the two coordinate functions with respect to \(X_c,Y_c,Z_c\).
\[
u=f_xX_c/Z_c+c_x,
\qquad
v=f_yY_c/Z_c+c_y
\]

\end{proof}

\begin{proposition}[Landmark Jacobian for the reprojection residual]
For the reprojection residual
\[
r_{ij}=z_{ij}-\pi(R_iX_j+t_i)
\]
The Jacobian with respect to the landmark \(X_j\) is
\[
\frac{\partial r_{ij}}{\partial X_j}
=
-
D\pi(Y_{ij})R_i,
\qquad
Y_{ij}=R_iX_j+t_i
\]
\end{proposition}

\begin{proof}
A perturbation \(X_j\mapsto X_j+\delta X_j\) gives
\[
Y_{ij}\mapsto Y_{ij}+R_i\delta X_j
\]
Hence
\[
\delta r_{ij}
=
-\,D\pi(Y_{ij})R_i\delta X_j
\]
\end{proof}

\begin{proposition}[Pose Jacobian under left \(SE(3)\) perturbation]
Let the camera pose be perturbed on the left by
\[
T_i\leftarrow \exp(\xi_i^\wedge)T_i
\qquad
\xi_i=
\begin{bmatrix}
\rho_i\\ \phi_i
\end{bmatrix}
\in\R^6.
\]
Then, to first order,
\[
\delta Y_{ij}
=
\rho_i-[Y_{ij}]_\times\phi_i
\]
Consequently,
\[
\frac{\partial r_{ij}}{\partial \xi_i}
=
-
D\pi(Y_{ij})
\begin{bmatrix}
I_3 & -[Y_{ij}]_\times
\end{bmatrix}
\]
\end{proposition}

\begin{proof}

Using
\[
\exp(\xi_i^\wedge)
\approx
\begin{bmatrix}
I+\phi_i^\wedge & \rho_i\\
0 & 1
\end{bmatrix}
\]
The perturbed camera-frame point is
\[
Y_{ij}'=(I+\phi_i^\wedge)Y_{ij}+\rho_i
=
Y_{ij}+\rho_i+\phi_i\times Y_{ij}
\]

Since
\[
\phi_i\times Y_{ij}=-[Y_{ij}]_\times\phi_i,
\]
We obtain
\[
\delta Y_{ij}
=
Y_{ij}'-Y_{ij}
=
\rho_i-[Y_{ij}]_\times\phi_i
\]
We know
\[
\delta r_{ij}
=
-\,D\pi(Y_{ij})\delta Y_{ij}
\]
So
\[
\delta r_{ij}
=
-
D \pi (Y_{ij})
\begin{bmatrix}
I_3 & -[Y_{ij}]_\times
\end{bmatrix}
\begin{bmatrix}
\rho_i\\ \phi_i
\end{bmatrix}
\]
\end{proof}


% ====
\section{Linearised Bundle Adjustment}
% ====

Bundle adjustment is nonlinear because both the camera projection and the pose parametrisation are nonlinear. We therefore try to solve a sequence of local linear least-squares problems.

\begin{definition}[Stacked residual and increment]
Let \( r \) denote the vector formed by stacking all residuals \( r_{ij} \)
\[
r(\theta)
\in\R^{2|\mathcal O|}
\]
Where \(\theta\) denotes all poses and landmarks. 

Let \(\delta\) denote the corresponding stacked local increment, consisting of pose increments in \(\R^6\) and landmark increments in \(\R^3\).
\end{definition}

\begin{proposition}[Local BA gives a linear least-squares problem]
At a current estimate \(\theta\), the residual admits the first-order approximation
\[
r(\theta+\delta)
=
r(\theta)+J\delta+O(\|\delta\|^2)
\]
Where \( J \) is the bundle adjustment Jacobian
\[
J=\frac{\partial r}{\partial \delta}
\]
Hence the local Gauss--Newton approximation to bundle adjustment is
\[
\min_\delta
\frac12
\|r+J\delta\|_2^2
\]
\end{proposition}

\begin{proof}
This is the first-order Taylor expansion of the stacked residual map \(r\) about the current estimate. Dropping terms of order \(O(\|\delta\|^2)\) gives the stated linear least-squares problem.
\end{proof}

\begin{proposition}[Normal equations for the Gauss--Newton step]
A stationary point of
\[
\min_\delta
\frac12
\|r+J\delta\|_2^2
\]
Satisfies
\[
J^\top J\delta=-J^\top r
\]
More generally, for a positive definite weight matrix \(W\), the weighted problem
\[
\min_\delta
\frac12
(r+J\delta)^\top W(r+J\delta)
\]
Has normal equations
\[
J^\top WJ\delta=-J^\top Wr
\]
\end{proposition}

\begin{proof}
Define
\[
F(\delta)=\frac12\|r+J\delta\|_2^2
\]
Then
\[
\nabla_\delta F
=
J^\top(r+J\delta)
\]
At a stationary point,
\[
J^\top(r+J\delta)=0
\]
Which gives
\[
J^\top J\delta=-J^\top r
\]

For the weighted case,
\[
F_W(\delta)=\frac12(r+J\delta)^\top W(r+J\delta)
\]
So
\[
\nabla_\delta F_W
=
J^\top W(r+J\delta)
\]
Setting the gradient to zero gives
\[
J^\top WJ\delta=-J^\top Wr
\]
\end{proof}

\begin{remark}
The matrix \(J^\top J\) is the Gauss--Newton approximation to the Hessian of the nonlinear least-squares objective, while \(J^\top r\) is the gradient. The central computational problem in bundle adjustment is therefore solving this linear system efficiently while exploiting its sparsity.
\end{remark}

% ====
\section{Sparse Structure}
% ====

The efficiency of bundle adjustment comes from the fact that each image measurement involves only one camera pose and one landmark. This gives the Jacobian, and hence the normal matrix, a highly sparse block structure.

\begin{definition}[Local Jacobian blocks]
For an observation \((i,j)\in\mathcal O\), write
\[
r_{ij}=r_{ij}(T_i,X_j)\in\R^2
\]
Let
\[
A_{ij}\coloneqq \frac{\partial r_{ij}}{\partial \xi_i}\in\R^{2\times 6},
\qquad
L_{ij}\coloneqq \frac{\partial r_{ij}}{\partial X_j}\in\R^{2\times 3}
\]
Where \(\xi_i\in\R^6\) is the local pose increment for camera \(i\), and \(X_j\in\R^3\) is landmark \(j\).
\end{definition}

\begin{proposition}[Each residual touches one pose and one landmark]
For
\[
r_{ij}=z_{ij}-\pi(R_iX_j+t_i),
\]
We have
\[
\frac{\partial r_{ij}}{\partial \xi_k}=0
\quad\text{for } k\neq i,
\qquad
\frac{\partial r_{ij}}{\partial X_\ell}=0
\quad\text{for } \ell\neq j.
\]
Thus each residual contributes non-zero Jacobian blocks only in the columns corresponding to camera \(i\) and landmark \(j\).
\end{proposition}

\begin{proof}
The residual \(r_{ij}\) is a function only of \(T_i\) and \(X_j\):
\[
r_{ij}=r_{ij}(T_i,X_j).
\]
It has no dependence on any other camera pose \(T_k\), \(k\neq i\), nor on any other landmark \(X_\ell\), \(\ell\neq j\). Therefore, the corresponding partial derivatives are zero.
\end{proof}

\begin{remark}
If all variables are ordered as
\[
\delta=
\begin{bmatrix}
\delta c\\
\delta p
\end{bmatrix}
\]
Where \(\delta c\in\R^{6M}\) stacks all camera increments and \(\delta p\in\R^{3N}\) stacks all landmark increments, then the Jacobian row block for \(r_{ij}\) has the form
\[
J_{ij}
=
\begin{bmatrix}
0 & \cdots & A_{ij} & \cdots & 0
&
0 & \cdots & L_{ij} & \cdots & 0
\end{bmatrix}
\]
\end{remark}

\begin{theorem}[BA normal matrix has camera--point block structure]
Consider the weighted linearised bundle adjustment problem
\[
\min_\delta
\frac12
\sum_{(i,j)\in\mathcal O}
(r_{ij}+A_{ij}\delta \xi_i+L_{ij}\delta X_j)^\top
\Omega_{ij}
(r_{ij}+A_{ij}\delta \xi_i+L_{ij}\delta X_j)
\]
Where \(\Omega_{ij}\in\R^{2\times 2}\) is positive definite. Then the normal equations have the block form
\[
\begin{bmatrix}
B & E\\
E^\top & C
\end{bmatrix}
\begin{bmatrix}
\delta c\\
\delta p
\end{bmatrix}
=
-
\begin{bmatrix}
g_c\\
g_p
\end{bmatrix}
\]
The non-zero blocks are
\[
B_{ii}
=
\sum_{j:(i,j)\in\mathcal O}
A_{ij}^\top \Omega_{ij}A_{ij}
\]
\[
C_{jj}
=
\sum_{i:(i,j)\in\mathcal O}
L_{ij}^\top \Omega_{ij}L_{ij}
\]
With
\[
E_{ij}
=
A_{ij}^\top \Omega_{ij}L_{ij}
\qquad \text{when } (i,j)\in\mathcal O
\]
The corresponding gradient blocks are
\[
(g_c)_i
=
\sum_{j:(i,j)\in\mathcal O}
A_{ij}^\top \Omega_{ij}r_{ij},
\qquad
(g_p)_j
=
\sum_{i:(i,j)\in\mathcal O}
L_{ij}^\top \Omega_{ij}r_{ij}
\]
\end{theorem}

\begin{proof}
The normal matrix is obtained by summing the per-observation contributions
\[
J_{ij}^\top \Omega_{ij} J_{ij}
\]
Since \(J_{ij}\) has non-zero blocks only in the columns for camera \(i\) and landmark \(j\), each observation contributes only the following four block terms:
\[
A_{ij}^\top\Omega_{ij}A_{ij},
\qquad
A_{ij}^\top\Omega_{ij}L_{ij},
\]
\[
L_{ij}^\top\Omega_{ij}A_{ij},
\qquad
L_{ij}^\top\Omega_{ij}L_{ij}
\]
Summing these terms over all observations gives the stated block matrix
\[
H=
J^\top\Omega J
=
\begin{bmatrix}
B & E\\
E^\top & C
\end{bmatrix}
\]
Similarly, the right-hand side is obtained from
\[
J^\top\Omega r
\]
Which gives the stated camera and landmark gradient blocks.
\end{proof}

\begin{proposition}[The landmark block \(C\) is block diagonal]
For pure reprojection bundle adjustment, the landmark--landmark block \(C\) satisfies
\[
C=\operatorname{blkdiag}(C_{11},\dots,C_{NN})
\]
Where
\[
C_{jj}
=
\sum_{i:(i,j)\in\mathcal O}
L_{ij}^\top \Omega_{ij}L_{ij}
\in\R^{3\times 3}
\]
In particular,
\[
C_{j\ell}=0
\qquad
\text{whenever } j\neq \ell
\]
\end{proposition}

\begin{proof}
A single residual \(r_{ij}\) depends on landmark \(X_j\) but on no other landmark. Therefore it cannot create a second-derivative or Gauss--Newton coupling between two distinct landmarks \(X_j\) and \(X_\ell\), \(j\neq \ell\). Equivalently, each per-observation contribution \(J_{ij}^\top\Omega_{ij}J_{ij}\) contains only the landmark--landmark block for landmark \(j\). 
\[
L_{ij}^\top\Omega_{ij}L_{ij}
\]
Summing over all observations preserves this block diagonal structure.
\end{proof}

\begin{remark}
The block diagonal form of \(C\) is the key algebraic fact behind Schur complement bundle adjustment. Since each block \(C_{jj}\) is only \(3\times3\), eliminating landmarks requires many small inversions rather than one large dense inversion.
\end{remark}

% ====
\section{Schur Complement}
% ====

Theorem: Landmark elimination by Schur complement
Proposition: Eliminating landmarks is cheap because C is block diagonal
Proposition: A landmark connects all cameras that observe it in the reduced system

% ====
\section{Gauge Freedom}
% ====

Theorem: Monocular BA is invariant under global Sim(3)
Corollary: The monocular BA Hessian has a 7-dimensional gauge nullspace
Proposition: Gauge fixing by anchoring poses, landmarks, or scale

% ====
\section{Damping and Robustification}
% ====

Proposition: LM damping makes the linear system positive definite
Proposition: LM interpolates between Gauss--Newton and gradient descent
Proposition: Robust losses induce iteratively reweighted least squares

% ====
\section{Practical BA Diagnostics}
% ====

Remark: Reprojection error distribution
Remark: Negative/near-zero depth
Remark: Poorly observed landmarks
Remark: Gauge drift
Remark: Outliers that survive RANSAC
Remark: Why local BA/windowed BA is used in SLAM


\end{document}